\chapter{Introduction}
\label{ch:introduction}

\section{Motivation}
\label{sec:motivation}

Almost every major securities and derivatives exchange organises trading
as a continuous double auction (CDA), implemented as a continuous limit
order book. Orders are processed one at a time in the order the exchange
received them, and whenever an incoming order crosses the opposite side of
the book a trade is executed immediately at the resting limit price. The
rule is simple, easy to explain, and a trader willing to pay the current
best price never has to wait.

The same rule has a far less benign structural consequence. Because
execution is serial and instantaneous, the value of a trading opportunity
accrues entirely to whoever reaches the matching engine first. When public
information changes the fair value of an asset, the stale quotes resting
in the book are mispriced for a short interval, and the first participant
to act either takes a stale quote or cancels it. Nothing about the
informational content of the event determines who wins; only relative
speed does. A mechanism designed to allocate goods according to
willingness to pay thus partly allocates them according to willingness to
invest in latency reduction \citep{budish2015hft}.

The competition is not costless. Microwave and hollow-core fibre links,
co-location and custom hardware absorb real resources without improving
the informational efficiency of prices in any first-order way, and for
participants outside the race the risk of being picked off is priced into
quotes: a market maker who expects to lose a fraction of its fills to
stale quotes must widen its spread to break even on the remainder.
\citet{aquilina2022quantifying} measure this cost directly from London
Stock Exchange message-level data.

The frequent batch auction (FBA) is the market design response proposed by
\citet{budish2015hft}. Instead of processing orders as they arrive, the
engine collects all orders arriving within a short interval $\tau$ and
clears them simultaneously at a single uniform price. Arrival order inside
the interval no longer matters, and all trades settle at the same price, so
a participant cannot obtain a better price by being earlier. Speed retains
value only across intervals, the marginal return to shaving another
nanosecond off a network path collapses to approximately zero, and
competition is redirected from latency to price.

The idea is not exotic: most equity exchanges already run uniform-price
call auctions at the open and close. What is new is the frequency --- an
auction every few hundred milliseconds or faster, so that the delay is
small for ordinary traders yet long enough to swamp the speed differences
that fund the arms race. In decentralised finance the argument arrives
from a different direction but converges on the same mechanism: settlement
is already discrete, ordering inside a block is economically contestable,
and uniform-price batch clearing removes the ordering rent otherwise
extracted by block producers \citep{gong2023fbadex,walther2021multitoken}.

\section{The Design Problem behind the Idea}
\label{sec:designproblem}

Stating that orders should be batched and cleared at a uniform price does
not yet specify a matching engine; it specifies a family of them. An
operator implementing one must decide how long the interval is and whether
its end is predictable, what participants see while orders accumulate,
whether all flow meets in one pool, and, when the two sides do not match
exactly at the clearing price, which rationing rule decides who is filled.
None of these questions has a single correct answer; each is a trade-off
that changes the strategic problem participants face. The theoretical part
of this thesis makes those trade-offs explicit and derives, from auction
theory and market microstructure theory, which behaviour each option should
be expected to induce.

\section{Problem Statement and Research Gap}
\label{sec:gap}

The literature on frequent batch auctions is largely about whether
batching is a good idea, not about how a batching engine is built.
\citet{budish2015hft} establish the equilibrium argument,
\citet{aquilina2022quantifying} quantify what the alternative costs,
\citet{zoican2021speed} and \citet{eibelshauser2022fba} study how
information and speed interact inside auctions, and
\citet{aldrich2020experiments} compare the two institutions in the
laboratory. Individual design dimensions appear in isolation: rationing
rules in the derivatives microstructure literature, disclosure in the
transparency literature \citep{madhavan1996transparency}, and clearing
formulations in the decentralised exchange literature
\citep{walther2021multitoken,ramsay2023speedex}. What is missing is a
consolidated treatment that decomposes the engine into its architectural
blocks, sets out the options for each, and traces each option to the
strategic behaviour it induces.

\paragraph{Field comparisons cannot hold order flow constant.}
The empirical literature comparing batch and continuous trading must
compare different venues, assets or time periods. \citet{lee2026taiwan}
exploit the Taiwan Stock Exchange's transition from intraday call auctions
to continuous trading, about as clean a natural experiment as this area
offers, yet the participant population, submission strategies and
information environment all adjust together with the mechanism. This is a
property of the object, not a shortcoming of the studies: an order stream
exists only under the mechanism that produced it, so a strictly
like-for-like comparison in which one and the same stream is executed once
continuously and once in batches cannot be obtained from field data.
Simulation with two interchangeable engines is the only design that holds
the flow fixed and varies the mechanism.

\section{Research Questions}
\label{sec:rqs}

The thesis is organised around two research questions, one for each part.

\begin{quote}
\textbf{RQ1 (theoretical).} How is a frequent batch auction engine
designed, and which design options exist for each of its architectural
components?
\end{quote}

RQ1 is answered by decomposing the engine into five design blocks --- the
batch interval, the information disclosure regime, order flow segregation,
the matching and rationing rule, and settlement and residual handling ---
enumerating the options within each, and deriving the strategic behaviour
each induces.

\begin{quote}
\textbf{RQ2 (empirical).} How do a continuous double auction and a
frequent batch auction differ when both engines are fed with the same
order flow?
\end{quote}

RQ2 is decomposed into three measurable dimensions, each operationalised by
a set of metrics defined in \Cref{app:formulas}:

\begin{itemize}
  \item \textbf{RQ2.1 Liquidity and transaction costs.} Which mechanism
  lets participants trade at a lower cost and with more depth available?
  \item \textbf{RQ2.2 Price discovery and market quality.} Which mechanism
  produces a more stable, less dispersed price series?
  \item \textbf{RQ2.3 Execution and order-arrival timing.} How much trades,
  how quickly and at what price relative to the submitted limit, and how
  does the batch mechanism reshape execution outcomes and arrival timing for
  the same unchanged order flow?
\end{itemize}

\section{Contribution}
\label{sec:contribution}

The thesis makes two contributions, one conceptual and one methodological.
It consolidates the FBA design space into five blocks and links each
available option to a theoretical prediction about induced behaviour. And
it implements two matching engines on a shared order and trade
representation, so that the mechanism is the only treatment variable and
the exact rationing outcome and the residual left unmatched after each
batch are directly observable rather than inferred, applying them to a
recently released Level~4 order book dataset from the Hyperliquid exchange
\citep{albers2026openbook} at nanosecond precision under identical
conditions with low confounding.

Both contributions are bounded by what a replay of fixed historical order
flow can show. The design decomposition is normative, not a report of
which option was tested; the empirical comparison isolates the
mechanism's mechanical effect on one dataset, not how participants would
have behaved had they known they were trading under it
(\Cref{sec:metricreliability,sec:limitations} set out exactly where that
line falls).
